\documentclass{article}

\usepackage{amsmath,amssymb}
\usepackage{xurl}
\usepackage[hidelinks]{hyperref}

% Shared commands for notes in docs/paper-gaps/.

\DeclareUnicodeCharacter{2080}{\ifmmode{}_{0}\else\textsubscript{0}\fi}
\DeclareUnicodeCharacter{2081}{\ifmmode{}_{1}\else\textsubscript{1}\fi}
\DeclareUnicodeCharacter{2082}{\ifmmode{}_{2}\else\textsubscript{2}\fi}
\DeclareUnicodeCharacter{2099}{\ifmmode{}_{n}\else\textsubscript{n}\fi}

\newcommand{\Lean}{\textsc{Lean}}
\ExplSyntaxOn
\NewDocumentCommand{\leanid}{m}{
  \ifmmode
    \textsf{\detokenize{#1}}
  \else
    \group_begin:
    \urlstyle{sf}
    \tl_set:Nn \l_tmpa_tl {#1}
    \tl_replace_all:Nnn \l_tmpa_tl {\_} {_}
    \exp_args:NV \path \l_tmpa_tl
    \group_end:
  \fi
}
\ExplSyntaxOff
\providecommand{\tr}{\operatorname{tr}}
\newcommand{\tnleanIssueBase}{https://github.com/LionSR/TNLean/issues/}
\newcommand{\tnleanPrBase}{https://github.com/LionSR/TNLean/pull/}
\newcommand{\tnleanDocBase}{https://sirui-lu.com/TNLean/docs/}
\newcommand{\ghissue}[1]{\href{\tnleanIssueBase#1}{GitHub issue \##1}}
\newcommand{\ghpr}[1]{\href{\tnleanPrBase#1}{PR \##1}}
\newcommand{\doclink}[1]{\href{\tnleanDocBase#1.html}{\path{#1}}}

% Dirac notation and the identity operator, as in blueprint/src/macros/common.tex.
% \providecommand keeps note-local definitions authoritative.
\usepackage{dsfont}
\providecommand{\ket}[1]{|#1\rangle}
\providecommand{\bra}[1]{\langle #1|}
\providecommand{\braket}[2]{\langle #1 | #2 \rangle}
\providecommand{\ketbra}[2]{|#1\rangle\!\langle #2|}
\providecommand{\Id}{\mathds{1}}
\providecommand{\one}{\mathds{1}}
\providecommand{\C}{\mathbb{C}}
\providecommand{\MN}[1]{M_{#1}(\C)}
\providecommand{\MD}{M_D(\C)}

% Verdict marker: \gapnote{<kind>}{<status>}. Kind is the policy
% classification (clarification, local-correction, scope-restriction,
% unfaithful, false-source, open-gap); status is open, wip (elimination
% underway), resolved, or historical. Machine-read by the site index;
% prints a verdict line.
\newcommand{\gapnote}[2]{\par\noindent\textbf{Verdict:} #1 (#2).\par}


\title{The Anomaly Class of the Two-Qubit $\mathbb{Z}_2\times\mathbb{Z}_2$
Tensors: From Detector Values to the Class}
\author{}
\date{2026-09-30}

\begin{document}
\maketitle
\gapnote{scope-restriction}{wip}

\section{The target statement}

Ref.~\cite{GarreRubio2023MPOSym}, subsection ``MPO algebras representing
$\mathbb{Z}_2\times\mathbb{Z}_2$'', states that
$H^3(\mathbb{Z}_2\times\mathbb{Z}_2,U(1))\cong\mathbb{Z}_2^3$ and lists
representatives of the classes. Its mixed type-II representative takes the
value $-1$ on the triples $(g_1g_2^i,g_1^jg_2,g_1^kg_2)$ with
$i,j,k\in\{0,1\}$ and $1$ elsewhere, where $g_1,g_2$ are the two generators.
Equivalently, in the coordinates $a=(a_1,a_2)$,
\begin{align}
    \omega_{\mathrm{II}}(a,b,c)
     & =(-1)^{a_1b_2c_2}.
    \label{eq:glm23z2z2_type_two_cocycle}
\end{align}
The source says that the values $\omega_g=\omega(g,g,g)$ in a normalized
gauge, at the three elements of order two, distinguish the classes. Its table
assigns to the type-II class the values
\begin{align}
    (\omega_a,\omega_b,\omega_{ab})
     & =(+1,+1,-1).
    \label{eq:glm23z2z2_type_two_detectors}
\end{align} The
periodic-boundary construction of the same reference represents a finite
group exactly by matrix product operators built from a three-cocycle. The
anomaly three-cocycle of an exact representation by normal tensors is
defined from its fusion tensors in Ref.~\cite{FrancoRubio2025MPUGauging}, in
the display preceding the three-cocycle equation. Its cohomology class does
not depend on the choice of fusion tensors. Neither source prints the
two-qubit tensors used here.

\section{The formalized statement}

The two-qubit tensors $M_e,M_x,M_y,M_{xy}$ of the condensation-defect
example are a construction of this development. Their periodic operators
satisfy $U_gU_h=\lambda(g,h)^LU_{gh}$ on a ring of $L$ sites, with
$\lambda(g,h)=(-1)^{g_1h_2}$, so they form only a projective family. They
equal the operators of the construction for the type-II cocycle followed by
the on-site operator $Z_2^{g_1}$, the Pauli $Z$ on the second qubit of each
site. Multiplying the letter of $M_g$ with input label $i$ by $(-1)^{g_1i_2}$
removes this factor. The resulting tensors $M'_g$ have the bond dimensions
of $M_g$, and their periodic operators are exactly those of the
construction, so they form an exact representation of
$\mathbb{Z}_2\times\mathbb{Z}_2$. Their doubled-index tensors are normal,
since the letters differ from those of $M_g$ by signs.

For this exact representation, let $\omega$ be the anomaly three-cochain of
any choice of fusion tensors. The formalization proves, for every
$g\in\mathbb{Z}_2\times\mathbb{Z}_2$,
\begin{align}
    \omega(g,e,g)\,\omega(g,g,g)
     & =\omega_{\mathrm{II}}(g,e,g)\,\omega_{\mathrm{II}}(g,g,g)
    =(-1)^{g_1g_2}.
    \label{eq:glm23z2z2_detector_identity}
\end{align}
The right side of \eqref{eq:glm23z2z2_detector_identity} is $+1$ at $(1,0)$
and $(0,1)$ and $-1$ at $(1,1)$. The left side is invariant under the gauge
action of the fusion tensors because $g^2=e$, and in a normalized gauge it is
$\omega(g,g,g)$. So these are the values \eqref{eq:glm23z2z2_type_two_detectors}
of the table. In particular the anomaly class of
the full group is nontrivial.\footnote{Formal declarations:
    \leanid{Z2Z2Condensation.kleinFamily\_isNormalRepresentation} in
    \doclink{TNLean/MPS/Examples/AnomalousCondensation/AnomalousCondensationZ2Z2Exact},
    and \leanid{Z2Z2Condensation.cyclicInvariant\_omega\_kleinFamily} and
    \leanid{Z2Z2Condensation.not\_isTrivialGaugeClass\_omega\_kleinFamily} in
    \doclink{TNLean/MPS/Examples/AnomalousCondensation/Z2Z2AnomalyClass}.}
The values are computed on the three subgroups of order two, by restricting
the family and its fusion tensors along the three embeddings of
$\mathbb{Z}_2$. The undressed projective family is given no anomaly class,
and the sources give it none.

\section{Resolution of detector separation}

The detector-separation step is now proved for arbitrary three-cocycles,
without a normalization hypothesis. Two three-cocycles have the same class
if and only if the three products
$\omega(g,e,g)\omega(g,g,g)$ agree at the nonidentity elements.
Every class has a unique representative $\omega_{pqr}$, and their parameters
add under multiplication. This gives
$H^3(\mathbb{Z}_2\times\mathbb{Z}_2,\mathbb{C}^{\times})\cong\mathbb{Z}_2^3$.
The proof normalizes an arbitrary cocycle by a coboundary, then uses one
square root in $\mathbb{C}^{\times}$ to fix six entries; the cocycle equation
and the three diagonal signs determine the remaining entries.\footnote{
    \leanid{TNLean.Algebra.ScalarThreeCochain.cohomologousTo\_iff\_klein\_three\_cyclicInvariants},
    \leanid{TNLean.Algebra.ScalarThreeCochain.existsUnique\_cohomologousTo\_kleinCocycleFamily},
    and \leanid{TNLean.Algebra.ScalarThreeCochain.kleinH3Equiv} in
    \doclink{TNLean/Algebra/KleinCocycleCompleteness}.}

Applied to \eqref{eq:glm23z2z2_detector_identity}, this proves
$[\omega]=[\omega_{\mathrm{II}}]$ for every choice of fusion tensors of the
exact dressed representation.\footnote{
    \leanid{Z2Z2Condensation.omega\_cohomologousTo\_kleinCocycle} in
    \doclink{TNLean/MPS/Examples/AnomalousCondensation/Z2Z2AnomalyClass}.}

\textbf{Verdict: detector separation resolved; the family comparison is still open.}
The conclusion concerns the exact dressed family $M'_g$. It does not assign
an anomaly class to the original odd-ring projective family, nor establish
transport to the even-ring or two-site-blocked original family. That family
comparison remains outside this result, as tracked by issue \#8041, so the
restriction stays live until it is settled.

\bibliographystyle{alpha}
\bibliography{references}

\end{document}
